
\documentclass[a4paper,11pt]{article}

\usepackage[utf8]{inputenc}
\usepackage[italian]{babel}
\usepackage{geometry}
\geometry{a4paper, margin=2.5cm}
\usepackage{hyperref}
\usepackage{titlesec}
\usepackage{enumitem}

\title{\textbf{Evoluzione delle Pipeline CI/CD verso l'AgentOps: Un'Analisi Empirica sui Repository di AI Agent}}
\author{Nome del Gruppo / Membri}
\date{\today}

\begin{document}

\maketitle

\begin{abstract}
L'automazione dei processi di sviluppo tramite strumenti di Continuous Integration e Continuous Deployment (CI/CD), come GitHub Actions (GA), è una pratica consolidata nell'ingegneria del software per garantire efficienza e qualità. Tuttavia, lo sviluppo di sistemi basati sull'Intelligenza Artificiale (AI) e, più recentemente, su Agenti Autonomi LLM, introduce nuove sfide legate al Technical Debt e alla sicurezza (AI Safety). Questo progetto mira a quantificare e analizzare l'evoluzione delle pipeline CI/CD in contesti "agent-enabled". Estraendo le cronologie dei workflow di GitHub Actions da una raccolta di repository open-source di agenti AI, confronteremo le metriche strutturali e le tipologie di modifiche con quelle dei progetti software tradizionali e AI-based. L'obiettivo è comprendere come le pratiche MLOps si stiano adattando per integrare pattern emergenti di AgentOps, come l'osservabilità, la valutazione delle traiettorie e i guardrail.
\end{abstract}

\section{Introduzione e Obiettivi (Cosa si vuole studiare)}
L'uso di GitHub Actions (GA) è diventato lo standard di fatto per l'automazione dei workflow nei repository GitHub, riducendo il carico di lavoro manuale per test, build e gestione delle dipendenze. La ricerca empirica ha ampiamente studiato i progetti software tradizionali, rivelando che i flussi di lavoro sono soggetti a modifiche continue, con una media di un aggiornamento ogni 159 giorni per file. In contesti tradizionali, l'ecosistema è altamente standardizzato, con il 98,48\% dei singoli step (task) che si affida ad "Actions" pubbliche prelevate dal Marketplace.

Tuttavia, i sistemi basati su Machine Learning (ML) presentano caratteristiche uniche che generano un "AI Technical Debt" (AITD). Tra i sintomi più comuni vi sono i "Pipeline Jungles", ovvero flussi di preparazione dati disordinati e difficili da mantenere, e un uso eccessivo di "Glue Code" per far comunicare pacchetti ML generici. Con l'avvento degli Agenti LLM, la natura non deterministica e autonoma dei modelli ha richiesto un ulteriore cambio di paradigma, portando all'emergere dell'AgentOps. Questo paradigma enfatizza l'osservabilità, necessaria per tracciare ragionamento, pianificazione e chiamate a tool esterni da parte degli agenti.

\textbf{Obiettivi del progetto:}
L'obiettivo primario di questo studio è analizzare empiricamente se e come le pipeline di CI/CD si differenzino in contesti "agent-enabled". Nello specifico, si intende:
\begin{itemize}
    \item Quantificare la complessità strutturale dei workflow (numero di job, step e dipendenze esterne) nei progetti AI agent open-source rispetto alle baseline del software tradizionale.
    \item Identificare la presenza di antipattern tipici dell'AITD nelle pipeline analizzate.
    \item Rilevare l'integrazione di pratiche specifiche dell'AgentOps, come l'inserimento di step di CI/CD dedicati al tracing logico, all'A/B testing dei prompt e ai guardrail di sicurezza.
\end{itemize}

\section{Raccolta Dati}
Per la raccolta dei dati, il progetto prenderà in esame i repository open-source di agenti AI elencati nella directory \textit{awesome-ai-agents}. Questa selezione garantirà un campione rappresentativo di sistemi agentici allo stato dell'arte. 

Piuttosto che analizzare solo l'ultima versione disponibile delle pipeline, si procederà con l'estrazione dell'intera evoluzione storica dei workflow. Per automatizzare questo processo, verranno utilizzati strumenti sviluppati dalla letteratura scientifica:
\begin{enumerate}
    \item \textbf{Estrazione dei file storici:} Utilizzeremo \texttt{gigawork} (Give me GitHub Actions Workflows), uno strumento open-source in Python progettato per estrarre l'intera cronologia dei commit che hanno interessato i file YAML situati nella cartella \texttt{.github/workflows}. Questo permetterà di raccogliere sia i workflow validi che i relativi metadati.
    \item \textbf{Estrazione delle differenze (Changesets):} Per comprendere la natura delle singole modifiche nel tempo, ci affideremo a \texttt{gawd} (GitHub Actions Workflow Differ). Questo tool Python consente di calcolare le differenze sintattiche tra una versione e la precedente del workflow, categorizzando automaticamente le modifiche in aggiunte, rimozioni, ridenominazioni o alterazioni di entità (come job e step).
\end{enumerate}

\section{Analisi dei Dati}
L'analisi dei dati si articolerà in una fase quantitativa (estrazione di metriche) e una fase qualitativa/testuale (ricerca di pattern), confrontando i risultati ottenuti con le baseline fornite dalla letteratura empirica sui progetti tradizionali.

\subsection{Analisi Quantitativa delle Metriche di Base}
Inizieremo mappando la struttura media delle pipeline CI/CD dei progetti Agent-based. Calcoleremo:
\begin{itemize}
    \item Il numero medio di workflow per repository, di job per workflow e di step per job. Valuteremo se questi numeri si discostano significativamente dalla media dei progetti tradizionali (che si attesta su 2,8 workflow, 1,8 job e 4,7 step).
    \item La frequenza di modifica (change rate) per determinare se la volatilità dei prompt o l'instabilità dei framework LLM portino a un tasso di aggiornamento dei workflow superiore alla media standard (aggiornati in media ogni 159 giorni).
    \item La percentuale di riutilizzo di Actions esterne dal Marketplace rispetto alla scrittura di script personalizzati diretti (direttiva \texttt{run}), per quantificare un'eventuale propensione al "Glue Code".
\end{itemize}

\subsection{Analisi Qualitativa: MLOps e AgentOps}
I cambiamenti estratti tramite \texttt{gawd} e i file YAML grezzi verranno analizzati tramite espressioni regolari e text mining per cercare evidenze di pratiche o "smell" specifici:
\begin{itemize}
    \item \textbf{Debito Tecnico AI (AITD):} Identificheremo la presenza di file di configurazione massicci o di \textit{Pipeline Jungles} intrecciati (script complessi di data-preparation nidificati direttamente all'interno dei file YAML).
    \item \textbf{Tracce di AgentOps:} Cercheremo chiamate specifiche a librerie di \textit{Observability} (es. Langfuse, Traceloop, AgentOps, Arize) che permettono il tracciamento degli \textit{span} di ragionamento e uso dei tool.
    \item \textbf{Trajectory Evaluation e Guardrails:} Verificheremo se gli step di test (tipicamente concentrati sulla "Task Configuration") si sono evoluti per includere valutazioni step-by-step o valutazioni delle traiettorie decisionali dell'agente, o se vi sono step specifici per l'attivazione e il testing dei \textit{guardrail} per la sicurezza AI.
\end{itemize}

L'esito di questa metodologia ibrida fornirà un quadro dettagliato di come l'avvento degli AI Agent stia ridefinendo lo standard dell'automazione software e rivelerà quanto le pipeline open-source attuali siano pronte a supportare il paradigma emergente dell'AgentOps.

\end{document}
